% ===============================================================
% INTRODUCTION GÉNÉRALE
% ===============================================================

\chapter*{Introduction générale}
\addcontentsline{toc}{chapter}{Introduction générale}

Le football est sans conteste le sport le plus suivi au monde, réunissant des milliards de passionnés à chaque grande compétition internationale. La Coupe du Monde 2026, organisée conjointement par le Maroc, les États-Unis et le Canada, s'annonce comme l'édition la plus grande de l'histoire avec 48 équipes participantes, offrant ainsi un contexte exceptionnel pour les applications de divertissement sportif numérique.

Dans ce contexte, les applications de fantasy sport connaissent une croissance exponentielle. Ces plateformes permettent aux fans de participer activement à leur passion en constituant des équipes virtuelles à partir de joueurs réels, dont les performances sont ensuite converties en points. Le marché mondial du fantasy sport était estimé à plus de 20 milliards de dollars en 2023 et devrait dépasser 48 milliards d'ici 2030.

\textbf{Problématique :} La majorité des plateformes de fantasy sport existantes (Fantasy Premier League, FotMob, Sorare) sont conçues pour les championnats nationaux réguliers et ne proposent pas d'expérience optimisée pour une compétition éphémère comme la Coupe du Monde. De plus, elles manquent souvent d'interactions sociales riches (batailles directes, chat en temps réel, prédictions communautaires) et d'une interface mobile moderne et fluide.

Ce projet a pour objectif général de concevoir et développer \textbf{WC Fantasy}, une application mobile complète dédiée à la Coupe du Monde 2026, permettant à des groupes d'amis de s'affronter dans des ligues privées ou publiques via un système de fantasy football complet et gamifié. Les objectifs spécifiques sont : (1) développer un backend sécurisé et performant avec Node.js/Express/PostgreSQL ; (2) concevoir une application mobile cross-platform avec Flutter ; (3) implémenter des fonctionnalités temps réel via Socket.io ; (4) intégrer un module de synchronisation automatique des scores via une API officielle.

La méthodologie adoptée est itérative : analyse des besoins, conception de l'architecture, développement par fonctionnalités, et validation continue. Le présent rapport est organisé en quatre chapitres : le Chapitre 1 présente le contexte et la problématique ; le Chapitre 2 analyse l'état de l'art ; le Chapitre 3 expose la conception du système ; le Chapitre 4 détaille le développement et les résultats obtenus.
